% ==============================================================================
% CAPITOLO 11: IL PRODOTTO NAZIONALE E LA SPESA AGGREGATA (MODELLO REDDITO-SPESA)
% ==============================================================================
\chapter{Il Prodotto Nazionale e la Spesa Aggregata}
\label{chap:prodotto_nazionale_spesa_aggregata}

\begin{tcolorbox}[colback=navyblue!5!white, colframe=navyblue, title=\textbf{Quadro Concettuale del Capitolo}]
Questo capitolo sviluppa il nucleo fondamentale della teoria macroeconomica keynesiana di breve periodo: il \textbf{modello reddito-spesa a due settori} (Famiglie e Imprese). Si formalizzano il principio della domanda effettiva, le funzioni di comportamento del consumo privato e del risparmio, la determinazione analitica e geometrica del reddito di equilibrio tramite la \textbf{Croce Keynesiana}, il meccanismo di aggiustamento basato sulle variazioni involontarie delle scorte, la derivazione del \textbf{moltiplicatore del reddito} mediante serie geometrica infinita e il celebre \textbf{Paradosso della Parsimonia}.
\end{tcolorbox}

\section{I Fondamenti della Teoria Keynesiana di Breve Periodo}

La macroeconomia moderna di breve periodo trae origine dai contributi di John Maynard Keynes (\textit{The General Theory of Employment, Interest and Money}, 1936), formulati in risposta alla Grande Depressione degli anni '30.

\subsection{Legge di Say vs Principio della Domanda Effettiva}
La visione classica e quella keynesiana differiscono radicalmente sulla causalità macroeconomica:

\begin{table}[htbp]
\centering
\small
\renewcommand{\arraystretch}{1.35}
\begin{tabularx}{\textwidth}{l X X}
\toprule
\textbf{Dimensione} & \textbf{Visione Classica / Neoclassica} & \textbf{Visione Keynesiana (Breve Periodo)} \\
\midrule
\textbf{Principio Cardine} & 
\textbf{Legge di Say}: \textit{"L'offerta crea la propria domanda"}. Ogni produzione genera redditi sufficienti ad acquistare l'output complessivo. & 
\textbf{Principio della Domanda Effettiva}: \textit{"La domanda determina la produzione"}. È la spesa programmata a determinare il livello del PIL. \\
\midrule
\textbf{Flessibilità Prezzi} & 
Prezzi e salari nominali sono \textbf{perfettamente flessibili} in tempo reale. & 
Prezzi e salari nominali sono \textbf{vischiosi / rigidi} (\textit{sticky prices}). \\
\midrule
\textbf{Equilibrio del Lavoro} & 
Il mercato del lavoro è sempre in \textbf{pieno impiego}; la disoccupazione è solo volontaria o frizionale. & 
Possono verificarsi equilibri di \textbf{sotto-occupazione} con disoccupazione involontaria e capacità produttiva in eccesso. \\
\midrule
\textbf{Ruolo dello Stato} & 
Neutralità (\textit{laissez-faire}); il mercato si autoregola spontaneamente. & 
Intervento attivo di stabilizzazione tramite politiche fiscali e monetarie per sostenere la domanda. \\
\bottomrule
\end{tabularx}
\caption{Confronto paradigmatico tra teoria classica e teoria keynesiana.}
\label{tab:say_vs_keynes}
\end{table}

\begin{legge}{Principio Keynesiano della Domanda Effettiva}{domanda_effettiva}
Nel breve periodo, in presenza di prezzi fissi e capacità produttiva non pienamente utilizzata, il livello di produzione aggregata ($Y$) e di occupazione ($L$) è determinato interamente dal livello della \textbf{Spesa Aggregata programmata} ($\SA$):
\begin{equation}
    Y = \SA
\end{equation}
Le imprese adeguano i volumi di produzione alla domanda effettivamente manifestata sul mercato, accumulando o decumulando scorte in caso di squilibrio temporaneo.
\end{legge}

\section{Le Componenti della Spesa Aggregata in un'Economia a Due Settori}

Nel modello base a due settori si assume un'economia chiusa agli scambi con l'estero ($NX = 0$) e senza intervento del settore pubblico ($G = 0, T = 0, TR = 0$).

La \textbf{Spesa Aggregata programmata} ($\SA$) si compone esclusivamente di consumi delle famiglie ($C$) e investimenti delle imprese ($I$):
\begin{equation}
    \SA = C + I
\end{equation}

\subsection{La Funzione del Consumo Keynesiana}

\begin{definizione}{Funzione del Consumo Keynesiana}{funzione_consumo}
La \textbf{funzione del consumo} descrive la relazione tra la spesa per consumi delle famiglie ($C$) e il reddito nazionale corrente ($Y$):
\begin{equation}
    C = C_0 + c Y \quad \text{con } C_0 > 0 \quad \text{e} \quad 0 < c < 1
\end{equation}
dove:
\begin{itemize}
    \item $C_0$ è il \textbf{Consumo Autonomo di sussistenza}: la spesa minima incomprimibile che le famiglie sosterrebbero anche con reddito nullo ($Y = 0$), finanziata attingendo a risparmi pregressi (\textit{desparmio}) o indebitandosi;
    \item $c = \PMgC = \dv{C}{Y}$ è la \textbf{Propensione Marginale al Consumo} (\textit{Marginal Propensity to Consume}): la frazione di ciascun euro aggiuntivo di reddito destinata all'acquisto di beni di consumo.
\end{itemize}
\end{definizione}

\begin{definizione}{Propensione Media al Consumo (PMeC)}{pmec}
La \textbf{Propensione Media al Consumo} ($\mathrm{PMeC}$) è il rapporto tra consumo totale e reddito complessivo:
\begin{equation}
    \mathrm{PMeC} = \frac{C}{Y} = \frac{C_0 + cY}{Y} = \frac{C_0}{Y} + c
\end{equation}
Poiché $C_0 > 0$, la $\mathrm{PMeC}$ è \textbf{strettamente decrescente} al crescere del reddito $Y$, ed è sempre superiore alla propensione marginale ($\mathrm{PMeC} > \PMgC$), convergendo asintoticamente a $c$ per $Y \to \infty$.
\end{definizione}

\subsection{La Funzione del Risparmio Complementare}

Poiché in assenza di Stato il reddito è integralmente ripartito tra consumi e risparmio ($Y \equiv C + S$), la funzione del risparmio privato ($S$) è data per differenza:
\begin{equation}
    S = Y - C = Y - (C_0 + cY) = -C_0 + (1 - c)Y
\end{equation}

\begin{definizione}{Funzione del Risparmio e Propensione Marginale al Risparmio}{funzione_risparmio}
\begin{equation}
    S = -C_0 + s Y \quad \text{con } s = \PMgS = 1 - c \in (0, 1)
\end{equation}
dove:
\begin{itemize}
    \item $-C_0$ è l'\textbf{intaccamento del risparmio} (risparmio negativo / desparmio) in corrispondenza di $Y = 0$;
    \item $s = \PMgS = \dv{S}{Y} = 1 - c$ è la \textbf{Propensione Marginale al Risparmio} (\textit{Marginal Propensity to Save}).
\end{itemize}
Per identità contabile fondamentale:
\begin{equation}
    \PMgC + \PMgS = c + s = 1 \iff \dv{C}{Y} + \dv{S}{Y} = 1
\end{equation}
\end{definizione}

\begin{definizione}{Reddito di Pareggio delle Famiglie (Break-Even Income)}{break_even}
Il \textbf{punto di pareggio} è il livello di reddito in corrispondenza del quale il consumo eguaglia esattamente il reddito ($C = Y$), cosicché il risparmio è nullo ($S = 0$):
\begin{equation}
    C = Y \iff C_0 + c Y_{\text{pareggio}} = Y_{\text{pareggio}} \implies Y_{\text{pareggio}} = \frac{C_0}{1 - c} = \frac{C_0}{s}
\end{equation}
\begin{itemize}
    \item Per $Y < Y_{\text{pareggio}}$: $C > Y \implies S < 0$ (le famiglie desparmiano);
    \item Per $Y > Y_{\text{pareggio}}$: $C < Y \implies S > 0$ (le famiglie accumulano risparmio netto).
\end{itemize}
\end{definizione}

\subsection{La Funzione dell'Investimento nel Modello Base}
Nel modello base a prezzi fissi, le decisioni di investimento delle imprese sono considerate \textbf{esogene}:
\begin{equation}
    I = I_0
\end{equation}
dove $I_0 > 0$ rappresenta la spesa autonoma per l'acquisto di nuovi beni capitali e impianti, determinata dallo stato delle aspettative a lungo termine degli imprenditori (\textit{animal spirits}).

\section{L'Equilibrio Macroeconomico e la Croce Keynesiana}

\subsection{Determinazione Algebrica del Reddito di Equilibrio}
La spesa aggregata programmata a due settori è:
\begin{equation}
    \SA = C + I = (C_0 + c Y) + I_0 = (C_0 + I_0) + c Y = A_0 + c Y
\end{equation}
dove $A_0 = C_0 + I_0$ è la \textbf{spesa autonoma totale} indipendente dal reddito.

\begin{modello}{Equilibrio Reddito-Spesa (Croce Keynesiana)}{croce_keynesiana}
La condizione di equilibrio macroeconomico richiede che la produzione programmata eguagli la spesa aggregata programmata:
\begin{equation}
    Y = \SA \iff Y = A_0 + c Y
\end{equation}
Risolvendo algebricamente per $Y$:
\begin{equation}
    Y - c Y = A_0 \implies Y(1 - c) = A_0 \implies Y^* = \frac{1}{1 - c} A_0 = \frac{1}{1 - c} (C_0 + I_0)
\end{equation}
\end{modello}

\begin{teorema}{Equivalenza tra Equilibrio $Y = \SA$ ed Equilibrio $S = I$}{equivalenza_s_i}
La condizione di equilibrio reddito-spesa $Y = \SA$ è matematicamente equivalente alla condizione di eguaglianza tra il risparmio programmato e gli investimenti programmati:
\begin{equation}
    Y = C + I \iff Y - C = I \iff S(Y^*) = I_0
\end{equation}
\end{teorema}

\begin{dimostrazione}[Equivalenza Risparmio-Investimento]
Sostituendo la funzione del risparmio $S = -C_0 + (1-c)Y$ nella condizione $S = I_0$:
\begin{equation}
    -C_0 + (1 - c)Y = I_0 \implies (1 - c)Y = C_0 + I_0 \implies Y^* = \frac{C_0 + I_0}{1 - c}
\end{equation}
che coincide esattamente con la soluzione ricavata da $Y = \SA$.
\end{dimostrazione}

\subsection{Rappresentazione Geometrica della Croce Keynesiana (TikZ G19)}

\begin{figure}[htbp]
\centering
\begin{tikzpicture}[scale=1.1, >=Stealth]
    % Assi cartesiani
    \draw[->, line width=1.2pt] (0,0) -- (10.5,0) node[right, font=\small\bfseries] {Reddito / Prodotto ($Y$)};
    \draw[->, line width=1.2pt] (0,0) -- (0,8.0) node[above, font=\small\bfseries] {Spesa Aggregata ($\SA$)};
    
    % Bisettrice a 45 gradi (Retta Y = SA)
    \draw[dashed, color=gray!80, line width=1.2pt] (0,0) -- (7.8,7.8) node[above right, font=\footnotesize\bfseries, color=gray!80] {Bisettrice $45^\circ$ ($Y = \SA$)};

    % Retta della Spesa Aggregata SA = A_0 + c Y
    \draw[color=navyblue, line width=2pt] (0, 2.4) -- (9.0, 6.9) node[right, font=\small\bfseries, color=navyblue] {$\SA = A_0 + c Y$};

    % Punto di Equilibrio E(Y*, SA*)
    \coordinate (E) at (4.8, 4.8);
    \fill[forestgreen] (E) circle (3.5pt) node[above left=3pt, font=\small\bfseries, color=forestgreen] {$E(Y^*, \SA^*)$};

    % Linee di proiezione dell'equilibrio
    \draw[dotted, line width=1pt, color=forestgreen] (4.8,0) node[below, font=\small\bfseries, color=forestgreen] {$Y^*$} -- (4.8,4.8);
    \draw[dotted, line width=1pt, color=forestgreen] (0,4.8) node[left, font=\small\bfseries, color=forestgreen] {$\SA^* = Y^*$} -- (4.8,4.8);

    % Intercetta Spesa Autonoma A_0
    \node[left, font=\small\bfseries, color=navyblue] at (0, 2.4) {$A_0 = C_0 + I_0$};

    % Caso Y_1 > Y* (Eccesso di Offerta / Accumulo Scorte)
    \coordinate (Y1_x) at (7.5, 0);
    \draw[dotted, line width=0.9pt, color=crimsonred] (7.5,0) node[below, font=\footnotesize\bfseries] {$Y_1$} -- (7.5, 7.5);
    \fill[crimsonred] (7.5, 6.15) circle (2.5pt); % punto su SA
    \fill[gray!80] (7.5, 7.5) circle (2.5pt); % punto su 45°
    \draw[<->, crimsonred, line width=1.4pt] (7.5, 6.25) -- (7.5, 7.4) node[midway, right, font=\scriptsize\bfseries] {$\Delta\mathrm{Scorte}_{\text{inv}} > 0$ (Accumulo)};
    \draw[->, crimsonred, line width=1.3pt] (7.2, 0.4) -- (5.4, 0.4) node[midway, above, font=\scriptsize\bfseries] {Taglio Produzione};

    % Caso Y_2 < Y* (Eccesso di Domanda / Decumulo Scorte)
    \coordinate (Y2_x) at (2.2, 0);
    \draw[dotted, line width=0.9pt, color=purpleaccent] (2.2,0) node[below, font=\footnotesize\bfseries] {$Y_2$} -- (2.2, 3.5);
    \fill[purpleaccent] (2.2, 3.5) circle (2.5pt); % punto su SA
    \fill[gray!80] (2.2, 2.2) circle (2.5pt); % punto su 45°
    \draw[<->, purpleaccent, line width=1.4pt] (2.2, 2.3) -- (2.2, 3.4) node[midway, left, font=\scriptsize\bfseries] {$\Delta\mathrm{Scorte}_{\text{inv}} < 0$ (Decumulo)};
    \draw[->, purpleaccent, line width=1.3pt] (2.4, 0.4) -- (4.2, 0.4) node[midway, above, font=\scriptsize\bfseries] {Aumento Produzione};

\end{tikzpicture}
\caption{La Croce Keynesiana (Diagramma a 45°). Determinazione del reddito di equilibrio $Y^*$ e dinamica di riequilibrio tramite la variazione non programmata delle scorte.}
\label{fig:croce_keynesiana}
\end{figure}

\section{Dinamica di Aggiustamento e Ruolo delle Scorte Involontarie}

Il meccanismo di trasmissione che garantisce la convergenza del sistema economico verso il punto di equilibrio $E$ è basato sul comportamento delle scorte aziendali:

\begin{definizione}{Variazione Non Programmata delle Scorte}{scorte_involontarie}
La \textbf{variazione involontaria delle scorte} ($\Delta \mathrm{Scorte}_{\text{inv}}$) è la differenza contabile tra la produzione corrente realizzata dalle imprese ($Y$) e la spesa programmata dagli acquirenti ($\SA$):
\begin{equation}
    \Delta \mathrm{Scorte}_{\text{inv}} = Y - \SA
\end{equation}
\end{definizione}

L'analisi dei tre regimi di mercato rivela:
\begin{enumerate}
    \item \textbf{Eccesso di Offerta ($Y > \SA$)}:
    \begin{itemize}
        \item La produzione realizzata eccede le vendite programmate;
        \item I beni invenduti rimangono nei magazzini $\implies \Delta \mathrm{Scorte}_{\text{inv}} > 0$ (\textbf{accumulo involontario di scorte});
        \item Segnale operativo: le imprese riducono gli ordini di produzione e l'impiego di fattori nel periodo successivo $\implies Y \downarrow$ fino a quando $Y = Y^*$.
    \end{itemize}
    
    \item \textbf{Eccesso di Domanda ($Y < \SA$)}:
    \begin{itemize}
        \item La domanda dei consumatori e delle imprese supera la produzione corrente;
        \item Le imprese evadono gli ordini intaccando le giacenze di magazzino $\implies \Delta \mathrm{Scorte}_{\text{inv}} < 0$ (\textbf{decumulo involontario di scorte});
        \item Segnale operativo: le imprese incrementano la produzione e assumono lavoratori $\implies Y \uparrow$ fino a raggiungere $Y^*$.
    \end{itemize}
    
    \item \textbf{Equilibrio Macroeconomico ($Y = \SA$)}:
    \begin{itemize}
        \item La produzione coincide esattamente con la spesa programmata;
        \item La variazione non programmata delle scorte è nulla ($\Delta \mathrm{Scorte}_{\text{inv}} = 0$);
        \item Il sistema è in stato stazionario di breve periodo.
    \end{itemize}
\end{enumerate}

\section{Il Moltiplicatore Keynesiano del Reddito}

\begin{teorema}{Il Moltiplicatore del Reddito}{moltiplicatore_teorema}
Dato il modello reddito-spesa $Y = \frac{1}{1 - c} A_0$, una variazione esogena della spesa autonoma ($\Delta A_0 = \Delta C_0 + \Delta I_0$) genera una variazione del reddito di equilibrio \textbf{più che proporzionale}:
\begin{equation}
    \Delta Y^* = \alpha_K \cdot \Delta A_0 \quad \text{con } \alpha_K = \frac{1}{1 - c} = \frac{1}{s} > 1
\end{equation}
dove $\alpha_K$ è il \textbf{moltiplicatore keynesiano}.
\end{teorema}

\begin{dimostrazione}[Derivazione del Moltiplicatore tramite Serie Geometrica Infinita]
Si consideri un incremento iniziale autonomo degli investimenti pari a $\Delta I_0$. L'impatto sul sistema economico si propaga attraverso una sequenza infinita di cicli di spesa e reddito:
\begin{enumerate}
    \item \textbf{Ciclo 1}: l'aumento degli investimenti genera una spesa immediata $\implies \Delta Y_1 = \Delta I_0$;
    \item \textbf{Ciclo 2}: l'aumento di produzione distribuisce un nuovo reddito di pari entità alle famiglie ($\Delta Y_1$), che ne consumano una quota $c$ $\implies \Delta C_1 = c \Delta I_0 \implies \Delta Y_2 = c \Delta I_0$;
    \item \textbf{Ciclo 3}: il nuovo consumo stimola ulteriore produzione e reddito, generando nuovi consumi $\implies \Delta C_2 = c \Delta Y_2 = c(c \Delta I_0) = c^2 \Delta I_0$;
    \item \textbf{Ciclo $k+1$}: al generico passo $k$, l'incremento di reddito indotto è $\Delta Y_{k+1} = c^k \Delta I_0$.
\end{enumerate}
La variazione complessiva cumulata del reddito è la somma di tutti gli infiniti incrementi:
\begin{equation}
    \Delta Y = \sum_{k=0}^{\infty} \Delta Y_{k+1} = \Delta I_0 + c \Delta I_0 + c^2 \Delta I_0 + c^3 \Delta I_0 + \dots = \Delta I_0 \sum_{k=0}^{\infty} c^k
\end{equation}
Poiché la propensione marginale al consumo è strettamente compresa tra 0 e 1 ($0 < c < 1$), la serie geometrica converge alla somma finita:
\begin{equation}
    \sum_{k=0}^{\infty} c^k = \frac{1}{1 - c}
\end{equation}
Pertanto:
\begin{equation}
    \Delta Y^* = \frac{1}{1 - c} \Delta I_0 = \alpha_K \Delta I_0
\end{equation}
\end{dimostrazione}

\subsection{Rappresentazione Grafica del Moltiplicatore (TikZ G20-G21)}

\begin{figure}[htbp]
\centering
\begin{tikzpicture}[scale=1.1, >=Stealth]
    % Assi
    \draw[->, line width=1.2pt] (0,0) -- (10.5,0) node[right, font=\small\bfseries] {Reddito ($Y$)};
    \draw[->, line width=1.2pt] (0,0) -- (0,8.0) node[above, font=\small\bfseries] {Spesa Aggregata ($\SA$)};
    
    % Bisettrice
    \draw[dashed, color=gray!80, line width=1.1pt] (0,0) -- (7.8,7.8) node[above right, font=\footnotesize\bfseries, color=gray!80] {$Y = \SA$};

    % Curva iniziale SA_0
    \draw[color=navyblue, line width=1.6pt] (0, 2.0) -- (8.5, 6.25) node[right, font=\footnotesize\bfseries, color=navyblue] {$\SA_0 = A_0 + c Y$};
    % Curva traslata SA_1 per aumento investimenti Delta I
    \draw[color=forestgreen, line width=2pt] (0, 3.4) -- (8.5, 7.65) node[right, font=\footnotesize\bfseries, color=forestgreen] {$\SA_1 = (A_0 + \Delta I) + c Y$};

    % Punti di equilibrio
    \coordinate (E0) at (4.0, 4.0);
    \coordinate (E1) at (6.8, 6.8);

    \fill[navyblue] (E0) circle (3pt) node[below right=2pt, font=\small\bfseries, color=navyblue] {$E_0(Y_0^*)$};
    \fill[forestgreen] (E1) circle (3.5pt) node[above left=2pt, font=\small\bfseries, color=forestgreen] {$E_1(Y_1^*)$};

    % Proiezioni
    \draw[dotted, line width=1pt, color=navyblue] (4.0,0) node[below, font=\small\bfseries] {$Y_0^*$} -- (E0);
    \draw[dotted, line width=1pt, color=forestgreen] (6.8,0) node[below, font=\small\bfseries] {$Y_1^*$} -- (E1);

    % Evidenziazione shift spesa Delta I
    \draw[<->, color=crimsonred, line width=1.5pt] (0.8, 2.4) -- (0.8, 3.8) node[midway, left, font=\footnotesize\bfseries] {$\Delta I_0$};

    % Evidenziazione variazione finale reddito Delta Y*
    \draw[<->, color=forestgreen, line width=1.5pt] (4.0, 0.5) -- (6.8, 0.5) node[midway, above, font=\footnotesize\bfseries] {$\Delta Y^* = \alpha_K \Delta I_0$};

\end{tikzpicture}
\caption{Rappresentazione geometrica del Moltiplicatore del Reddito. L'incremento verticale della spesa autonoma $\Delta I_0$ genera un incremento orizzontale del PIL $\Delta Y^*$ strettamente maggiore ($\Delta Y^* > \Delta I_0$).}
\label{fig:moltiplicatore_keynesiano}
\end{figure}

\section{Il Paradosso della Parsimonia}

\begin{modello}{Il Paradosso della Parsimonia}{paradosso_parsimonia}
Il \textbf{Paradosso della Parsimonia} (\textit{Paradox of Thrift}) afferma che il tentativo collettivo della società di aumentare il risparmio aggregato (attraverso un incremento della propensione marginale al risparmio $s$ o una riduzione del consumo autonomo $C_0$) produce una contrazione del reddito nazionale di equilibrio, lasciando \textbf{completamente inalterato il livello aggregato finale di risparmio effettivo}:
\begin{equation}
    S^* = I_0 = \text{costante}
\end{equation}
\end{modello}

\begin{dimostrazione}[Dimostrazione Analitica del Paradosso della Parsimonia]
Si consideri il livello di risparmio privato all'equilibrio macroeconomico:
\begin{equation}
    S(Y^*) = -C_0 + s Y^*
\end{equation}
Sostituendo il valore di equilibrio del reddito $Y^* = \frac{C_0 + I_0}{s}$:
\begin{equation}
    S^* = -C_0 + s \left[ \frac{C_0 + I_0}{s} \right] = -C_0 + (C_0 + I_0) = I_0
\end{equation}
Poiché gli investimenti $I_0$ sono costanti ed esogeni, il risparmio finale $S^*$ dipende unicamente da $I_0$.

Calcolando la derivata del reddito di equilibrio rispetto alla propensione al risparmio $s$:
\begin{equation}
    \dv{Y^*}{s} = -\frac{C_0 + I_0}{s^2} < 0
\end{equation}
Un aumento della virtù del risparmio da parte delle famiglie riduce la spesa complessiva, provocando una recessione che distrugge esattamente la base di reddito da cui il risparmio doveva essere generato.
\end{dimostrazione}

\begin{figure}[htbp]
\centering
\begin{tikzpicture}[scale=1.1, >=Stealth]
    % Assi
    \draw[->, line width=1.2pt] (0,0) -- (10.5,0) node[right, font=\small\bfseries] {Reddito ($Y$)};
    \draw[->, line width=1.2pt] (0,-2.5) -- (0,5.5) node[above, font=\small\bfseries] {Risparmio ($S$), Investimenti ($I$)};
    
    % Retta Investimenti costante I = I_0
    \draw[color=forestgreen, line width=1.8pt] (0, 2.5) -- (9.5, 2.5) node[right, font=\small\bfseries, color=forestgreen] {$I = I_0$};

    % Curva del risparmio iniziale S_0 = -C_0 + s_0 Y
    \draw[color=navyblue, line width=1.6pt] (0, -1.8) -- (9.0, 4.2) node[right, font=\small\bfseries, color=navyblue] {$S_0 = -C_0 + s_0 Y$};
    % Curva del risparmio con maggiore propensione al risparmio S_1 (ruotata verso l'alto)
    \draw[color=crimsonred, line width=2pt] (0, -1.8) -- (7.5, 5.0) node[above right, font=\small\bfseries, color=crimsonred] {$S_1 = -C_0 + s_1 Y \quad (s_1 > s_0)$};

    % Intercetta -C_0
    \node[left, font=\footnotesize\bfseries, color=navyblue] at (0, -1.8) {$-C_0$};

    % Punti di equilibrio
    \coordinate (E0) at (6.45, 2.5);
    \coordinate (E1) at (4.75, 2.5);

    \fill[navyblue] (E0) circle (3pt) node[above right, font=\small\bfseries, color=navyblue] {$E_0$};
    \fill[crimsonred] (E1) circle (3.5pt) node[above left, font=\small\bfseries, color=crimsonred] {$E_1$};

    % Proiezioni del reddito
    \draw[dotted, line width=1pt, color=navyblue] (6.45, 0) node[below, font=\small\bfseries] {$Y_0^*$} -- (E0);
    \draw[dotted, line width=1pt, color=crimsonred] (4.75, 0) node[below, font=\small\bfseries] {$Y_1^*$} -- (E1);

    % Freccia contrazione del reddito
    \draw[->, crimsonred, line width=1.5pt] (6.3, 0.4) -- (4.9, 0.4) node[midway, above, font=\footnotesize\bfseries] {$\Delta Y^* < 0$};

\end{tikzpicture}
\caption{Rappresentazione geometrica del Paradosso della Parsimonia. L'aumento della propensione al risparmio ($s_1 > s_0$) ruota la curva del risparmio verso l'alto, riducendo il PIL di equilibrio da $Y_0^*$ a $Y_1^*$, mentre il livello finale del risparmio rimane invariato a $I_0$.}
\label{fig:paradosso_parsimonia}
\end{figure}

\section{Metodi Risolutivi ed Esercizi Guidati Svolti}

\begin{metodo}[Algoritmo per la Risoluzione del Modello Reddito-Spesa a 2 Settori]
Per determinare l'equilibrio macroeconomico e quantificare l'effetto di shock esogeni:
\begin{enumerate}
    \item Identificare i parametri della funzione del consumo ($C_0, c$) e degli investimenti ($I_0$);
    \item Calcolare la spesa autonoma totale: $A_0 = C_0 + I_0$;
    \item Calcolare il moltiplicatore keynesiano: $\alpha_K = \frac{1}{1 - c}$;
    \item Calcolare il reddito di equilibrio: $Y^* = \alpha_K A_0$;
    \item Calcolare il consumo di equilibrio $C^* = C_0 + c Y^*$ e verificare che $S^* = Y^* - C^* = I_0$;
    \item In presenza di uno shock $\Delta I_0$, calcolare la variazione del PIL: $\Delta Y^* = \alpha_K \Delta I_0$ e il nuovo reddito $Y_1^* = Y_0^* + \Delta Y^*$.
\end{enumerate}
\end{metodo}

\begin{esame}[Tranelli Frequenti sul Modello Reddito-Spesa]
\begin{itemize}
    \item \textbf{Confusione tra Propensione Marginale e Media}: ricordare che $\PMgC = c$ è la pendenza della retta del consumo (costante nel modello lineare), mentre $\mathrm{PMeC} = C_0/Y + c$ decresce con il reddito.
    \item \textbf{Scorte non programmate fuori equilibrio}: se all'esame viene chiesto cosa accade per un dato livello di produzione $Y_a \neq Y^*$, calcolare $\SA(Y_a) = A_0 + c Y_a$ e determinare $\Delta \mathrm{Scorte} = Y_a - \SA(Y_a)$. Se $\Delta \mathrm{Scorte} > 0 \implies$ accumulo involontario; se $\Delta \mathrm{Scorte} < 0 \implies$ decumulo.
    \item \textbf{Variazione percentuale vs Variazione assoluta}: se l'investimento aumenta del $10\%$, la variazione assoluta è $\Delta I = 0,10 \times I_0$, e la variazione del PIL è $\Delta Y^* = \alpha_K \Delta I$.
\end{itemize}
\end{esame}

\begin{esercizio}[Calcolo Completo dell'Equilibrio, Scorte e Moltiplicatore]
In un'economia a due settori sono note le seguenti equazioni di comportamento:
\begin{align*}
    C &= 100 + 0{,}80 Y \\
    I &= 60
\end{align*}

\textbf{Quesiti}:
\begin{enumerate}
    \item Determinare la propensione marginale al risparmio, il moltiplicatore keynesiano e il reddito di equilibrio $Y^*$;
    \item Calcolare il consumo $C^*$, il risparmio $S^*$ e verificare la coerenza con gli investimenti;
    \item Valutare cosa accade nel sistema se la produzione corrente è $Y = 900$;
    \item Determinare la variazione del reddito di equilibrio se le imprese aumentano gli investimenti di $\Delta I = 25\,\text{\euro}$.
\end{enumerate}

\textbf{Svolgimento}:
\begin{enumerate}
    \item \textbf{Propensioni e Moltiplicatore}:
    \begin{align*}
        c &= \PMgC = 0{,}80 \implies s = \PMgS = 1 - 0{,}80 = 0{,}20 \\
        \alpha_K &= \frac{1}{1 - c} = \frac{1}{1 - 0{,}80} = \frac{1}{0{,}20} = 5
    \end{align*}
    La spesa autonoma totale è $A_0 = C_0 + I_0 = 100 + 60 = 160\,\text{\euro}$.\\
    Il reddito di equilibrio è:
    \begin{equation*}
        Y^* = \alpha_K \cdot A_0 = 5 \times 160 = 800\,\text{\euro}
    \end{equation*}
    
    \item \textbf{Consumo e Risparmio di Equilibrio}:
    \begin{align*}
        C^* &= 100 + 0{,}80 \times 800 = 100 + 640 = 740\,\text{\euro} \\
        S^* &= Y^* - C^* = 800 - 740 = 60\,\text{\euro}
    \end{align*}
    Si riscontra l'uguaglianza $S^* = I_0 = 60\,\text{\euro}$.
    
    \item \textbf{Analisi Fuori Equilibrio per $Y = 900$}:
    La spesa programmata per $Y = 900$ è:
    \begin{equation*}
        \SA(900) = 160 + 0{,}80 \times 900 = 160 + 720 = 880\,\text{\euro}
    \end{equation*}
    La variazione non programmata delle scorte è:
    \begin{equation*}
        \Delta \mathrm{Scorte}_{\text{inv}} = Y - \SA = 900 - 880 = +20\,\text{\euro}
    \end{equation*}
    Poiché $\Delta \mathrm{Scorte} > 0$, si verifica un \textbf{accumulo involontario di scorte}; le imprese riducono la produzione finché il PIL non torna al livello di equilibrio $Y^* = 800$.
    
    \item \textbf{Effetto dello Shock $\Delta I = 25\,\text{\euro}$}:
    \begin{align*}
        \Delta Y^* &= \alpha_K \cdot \Delta I = 5 \times 25 = +125\,\text{\euro} \\
        Y_1^* &= Y_0^* + \Delta Y^* = 800 + 125 = 925\,\text{\euro}
    \end{align*}
\end{enumerate}
\end{esercizio}
